\documentclass[12pt, a4paper]{article} % Classe: article, report o book [8]

% --- Lingua e Codifica (Supporto Italiano) ---
\usepackage[utf8]{inputenc}  % Codifica input
\usepackage[T1]{fontenc}      % Codifica font
\usepackage[italian]{babel}   % Lingua italiana [2]

% --- Impaginazione e Margini ---
\usepackage[margin=2.5cm]{geometry} % Imposta margini (2.5cm su tutti i lati)

% --- Pacchetti Scientifici/Matematici ---
\usepackage{amsmath, amssymb, amsfonts, amsthm} % Matematica avanzata [2]
\usepackage{graphicx} % Per inserire immagini
\usepackage{hyperref} % Per collegamenti ipertestuali
\usepackage{booktabs} % Per tabelle professionali

% --- Formattazione Testo e Altro ---
\usepackage[document]{ragged2e}
\usepackage{xcolor}    % Per usare i colori 


\title{Nicholas Fadda}
\date{23-01-2026}


\begin{document}

\maketitle

\section{Ruolo}
Nicholas Fadda è un Agente esterno Cepi che lavora in Geotech con esperienza di
circa un anno di collaborazione con l'azienda e si occupa di :
\begin{itemize}
      \item Creazione delle offerte Access per conto di clienti della sua zona di
            competenza.
      \item Gestione delle relazioni con i clienti.
      \item Supporto post-vendita e assistenza ai clienti.
\end{itemize}
Nicholas oltre ad Access utilizza anche intrenemnate alla sua azienda il software Pipedrive per la gestione dei clienti e delle fasi di avanzamento.
Nicholas suggerisce che internamente a cepi si utilizzi Hubspot per la gestione dei clienti e dei lead commerciali.

\section{Attuale uso di Access}
Nicholas utilizza Access per:
\begin{itemize}
      \item Creare le offerte per i clienti.
      \item Creazione delle offerte di assistenza e ampliamenti.
      \item Aggiunta clienti e gestione anagrafiche.
      \item Trasfomazione delle offerte in ordini se accettate.
\end{itemize}

\subsection{Attuale vista dati in Access}
Nicholas all'interno dell'attuale programma ha la possibilità di vedere:
\begin{itemize}
      \item Le offerte create da lui o dal suo referente (Michele Foietta).
      \item Le richieste di assistenza (RA) create da lui.
      \item L'anagrafica clienti completa.
      \item L'elenco articoli completo.
      \item Le offerte e RA degli altri agenti non sono visibili.
\end{itemize}

\section{Problemi attuali del sistema}
Nicholas ha evidenziato i seguenti problemi con l'attuale sistema Access:
\begin{itemize}
      \item \textbf{Instabilità e lentezza:} Il sistema Access attuale soffre di blocchi frequenti che costringono al riavvio forzato,
            specialmente durante operazioni di stampa o salvataggio.
      \item \textbf{Dipendenza da VPN:} La connessione cade spesso in mobilità (hotspot),
            causando il crash della connessione e l'impossibilità di operare davanti al cliente in real-time.
      \item \textbf{Ricerca limitata:} Impossibilità di cercare articoli per codice; la ricerca avviene solo per descrizione che però poco intuitive.
            Se abbiamo il codice gestionale esatto non esiste una corrispondenza nell'attuale programma.
            Per ovviare questo problema si utilizzano anche componenti erronei per dare un prezzo.
      \item \textbf{Mancanza di storico:} Nicholas non ha accesso allo storico globale o a offerte precedenti al 2025. Ciò limita la capacità di replicare offerte passate.
      \item \textbf{Mono-finestra:} Impossibilità di aprire più offerte contemporaneamente per confrontare o copiare dati.

\end{itemize}

\section{Requisiti funzionali per il nuovo software}

\subsection{Gestione offerte e workflow}
La gestione delle offerte secondo Nicholas deve prevedere:
\begin{itemize}
      \item \textbf{Modulo registrazione cliente:} Attualmente Nicholas usa un modulo creato da lui e Foietta dove scrive i principali dati di un cliente nel primo incontro in modo da essere veloce, ma efficiente. Solo succesivamnte crea un'offerta.
            Ha espresso la possibilità di impementare un modulo facilmnete editabile per eviatre di creare contrati subito verborsi e lunghi.
      \item \textbf{Versioning delle offerte:} Necessità di gestire le revisioni \textit{(es. Rev. A, Rev. B)} mantenendo lo stesso numero di offerta base,
            invece di duplicare il file creando nuovi numeri di offerta. Possibilità quindi di promuovere ad ordine solo una delle versioni.
      \item \textbf{Struttura per capitoli:} Mantenere e migliorare la suddivisione dell'offerta in capitoli logici \textit{(es. Silos, Trasporto, Automazione)}
            per chiarezza verso il cliente.
      \item \textbf{Template per tipo impianto:} Possibilità di creare modelli di offerta predefiniti per tipologie comuni di impianti
            (es. Pan di spagna , Biscotti , ecc \dots) per velocizzare la creazione e verificare che si è inserito l'occorrente.
      \item \textbf{Importazione intelligente:} Funzionalità per importare righe da altri documenti/offerte esistenti, selezionando quali voci copiare.
      \item \textbf{Calcolo installazione:} Modulo per il calcolo automatico dei costi di montaggio basato su variabili
            (N. tecnici $\times$ N. giorni $\times$ Tariffa giornaliera).
      \item \textbf{Gestione sconti:} L'attuale sistema di sconti è da mantenere:
            \begin{itemize}
                  \item Sconto sul singolo articolo.
                  \item Sconto sui capitoli.
                  \item Sconto globale sul materiale.
                  \item Sconto totale comprensivo di materiale e montaggio.
                  \item Possibilità di imporre un prezzo.
            \end{itemize}
\end{itemize}

\subsection{Catalogo e configurazione prodotto}
Per una corretta gestione del catalogo e della configurazione prodotto, il sistema deve prevedere:
\begin{itemize}
      \item \textbf{Ricerca avanzata:} Ricerca per Codice, Descrizione, Tag e Argomento (Category-based).
      \item \textbf{Configurazione guidata (Guided Selling):} Selezionando un componente principale \textit{(es. Silos in Trevira)},
            il sistema deve proporre/flaggare automaticamente gli accessori compatibili \textit{(sonde, valvole)} ed escludere quelli errati.
      \item \textbf{Standardizzazione codici:} Allineamento tra i codici commerciali e quelli tecnici \textit{(AutoCAD/Ufficio Tecnico)}
            per evitare discrepanze in fase di ordine e facilitare l'inserimento dei codici precisi.
      \item \textbf{Descrizioni componenti} Prestare attenzione alle descrizioni degli articoli per renderle più chiare e comprensibili.
            Attualmente sono presenti delle descrizioni poco chiare.
      \item \textbf{Standardizzazione:} Richiesta una standardizzazione dei modelli più usati per facilitare le offerte e renderle più simili tra loro.
\end{itemize}

\subsection{Gestione cliente e CRM (On-Site)}
Per una gestione efficace del cliente e delle anagrafiche, il sistema deve prevedere:
\begin{itemize}
      \item \textbf{Note tecniche permanenti:} Campi dedicati nell'anagrafica impianto per salvare specifiche strutturali
            (es. altezza soffitto, tipologia muri sandwich/cemento, accessibilità scarico).
      \item \textbf{Termini commerciali default:} Associazione automatica delle condizioni di pagamento e consegna specifiche per il cliente selezionato.
      \item \textbf{Origine contatto:} Campo per tracciare la fonte del lead (es. Fiera Sigep, Referral San Cassiano, ecc.).
\end{itemize}

\subsection{Integrazione e visualizzazione}
Per migliorare l'integrazione con altri sistemi e la visualizzazione delle informazioni, il sistema deve prevedere:
\begin{itemize}
      \item \textbf{Repository disegni/media:} Integrazione (es. OneDrive/Cartelle Server) per mostrare al cliente esempi di realizzazioni,
            PDF e DWG (senza dati sensibili di altri clienti) direttamente dall'interfaccia.
      \item \textbf{Tracking commessa:} Visualizzazione semplificata dello stato avanzamento commessa
            (Project Management $\rightarrow$ Produzione $\rightarrow$ Spedizione) accessibile ai commerciali.
\end{itemize}

\section{Requisiti tecnici e UX}
Nicola suggerisce i seguenti requisiti tecnici e di esperienza utente per il nuovo software:
\begin{itemize}
      \item \textbf{Modalità offline/sync:} Il software deve permettere di lavorare in locale (offline) durante le visite dai clienti
            e sincronizzarsi automaticamente con il database centrale una volta ristabilita la connessione VPN.
      \item \textbf{Multi-finestra:} Possibilità di tenere aperte più schede o finestre di offerte diverse contemporaneamente.
            Utile la possibilita di un sistema multi-documento \textit{Es: AutoCAD}.
      \item \textbf{Interfaccia mobile/tablet:} Ottimizzazione per l'uso su tablet (iPad) per la compilazione di dati preliminari e note durante il sopralluogo o in fiera.
\end{itemize}

\section{Suggerimenti futuri}
Per miglioramenti futuri anche esterni dal configuratore, Nicholas propone:
\begin{itemize}
      \item \textbf{QR Code/NFC per ricambi:} Generazione di QR Code sugli accessori per facilitare la
            richiesta di ricambi da parte del cliente, con notifica diretta al commerciale di riferimento.
      \item \textbf{Sconti su singolo articolo:} Possibilità di applicare sconti specifici su righe (es. materiale usato/da fiera)
            indipendentemente dallo sconto globale.
      \item \textbf{Unificazione blocchi schema commerciali/tecnici} : Unificare i blocchi AutoCAD di disegno per gli schemi tra gli schemi commerciali e tecnici
            per evitare discrepanze e facilitare la comunicazione tra reparti.
            Attualmente si usano blocchi di disegno differenti con simbologie diverse.
      \item \textbf{Integrazione con CRM esterno:} Collegamento diretto con software CRM esterni (es. Hubspot, Pipedrive)
            per sincronizzare automaticamente i dati cliente e le offerte create.
      \item \textbf{Gestione ricambistica:} Implementare delle query che recuperano i possibili ricambi in base al venduto precedente.
      \item \textbf{Avanzamneto ordine:} Implementare un sistema di avanzamento ordine che permetta di seguire lo stato di avanzamento
            dell'ordine in tutte le sue fasi (produzione, spedizione, installazione) per dare modo al commerciale di tenere
            informato il cliente senza interpellare altri reparti. Basta una casella di spunta per ogni fase.
\end{itemize}

\section{Tempi operativi}
Nicholas ha detto che per un nuovo assunto il tempo di complizione con
l'attuale sistema può arrivare anche ad un ora per singola offerta. Con un
sistema più intuitivo e veloce come quello descritto potrebbe migliorare
notevolmente i tempi di completamento delle offerte.

\end{document}